% ==========================================
% BAB VII PENUTUP
% ==========================================
\cleardoublepage
\chapter{PENUTUP}
\label{chap:penutup}

Bab ini menyajikan kesimpulan dari pelaksanaan tugas akhir berdasarkan tujuan yang telah ditetapkan, serta saran untuk pengembangan sistem di masa mendatang.

\section{Kesimpulan}

Berdasarkan hasil perancangan, implementasi, dan evaluasi yang telah dilakukan, diperoleh kesimpulan sebagai berikut.

\begin{enumerate}
  \item Arsitektur sistem \textit{Business Intelligence} berbasis \textit{chatbot} telah berhasil dirancang dan diimplementasikan menggunakan arsitektur tiga lapis yang terdiri dari \textit{frontend} React.js, \textit{backend} Node.js/Express, dan layanan \textit{machine learning} FastAPI/PyTorch. Sistem mampu memproses permintaan analisis data dan memberikan \textit{insight} berbasis data kepada pengguna internal perusahaan secara otomatis melalui antarmuka percakapan berbasis web.

  \item Modul pemrosesan bahasa alami telah diimplementasikan menggunakan pendekatan pencocokan pola kata kunci berbobot (\textit{weighted keyword matching}) yang mendukung tujuh \textit{intent} utama dan delapan fokus analisis. Pendekatan ini memungkinkan pengguna internal mengajukan pertanyaan dan permintaan analisis bisnis menggunakan bahasa Indonesia alami melalui antarmuka \textit{chatbot} dengan latensi klasifikasi kurang dari 1 milidetik.

  \item Mekanisme \textit{rule-based query} telah diimplementasikan melalui mesin aturan (\textit{rule engine}) yang mendukung lebih dari 30 aturan analitik, mencakup lima domain basis data bisnis (penjualan, pendapatan, target, proses data pelanggan, dan \textit{churn}). Mesin aturan dilengkapi fitur injeksi filter geografis secara dinamis dan mekanisme \textit{caching} untuk menghasilkan respons yang cepat dan konsisten.

  \item Model peramalan deret waktu berbasis \textit{deep learning} telah diintegrasikan ke dalam sistem menggunakan arsitektur GRU (\textit{Gated Recurrent Unit}) dan LSTM (\textit{Long Short-Term Memory}) dengan \textit{pipeline} pelatihan yang mencakup penyesuaian otomatis ukuran jendela (\textit{auto-tune window size}), validasi \textit{walk-forward}, penghentian dini (\textit{early stopping}), dan mekanisme pembandingan otomatis terhadap model sebelumnya. Berdasarkan evaluasi komparatif terhadap empat model menggunakan sistem penilaian KPI tiga kategori (LULUS, Lulus Dengan Catatan, GAGAL), model GRU memberikan kinerja paling \textit{robust} untuk metrik \textit{avg\_install\_days} dengan lima wilayah berstatus LULUS dan satu wilayah (REG-4) berstatus Lulus Dengan Catatan, tanpa ada wilayah yang berstatus GAGAL. Model ARIMA kompetitif pada deret waktu stabil namun gagal pada wilayah bervolatilitas tinggi, sementara model \textit{Prophet} menghasilkan sMAPE di atas $100\%$ pada seluruh wilayah karena asumsi dekomposisi aditif tidak sesuai dengan karakteristik data operasional.

  \item Modul pembangkitan bahasa alami berbasis templat (\textit{template-based Natural Language Generation}) telah diintegrasikan untuk menghasilkan narasi dalam bahasa Indonesia dari data terstruktur hasil kueri \textit{Enterprise Data Warehouse} Oracle. Akurasi faktual narasi mencapai 100\% karena seluruh data numerik berasal langsung dari hasil kueri basis data.

  \item Antarmuka percakapan berbasis web telah dirancang dan diimplementasikan menggunakan React.js dengan dukungan \textit{render Markdown} untuk menampilkan tabel dan format teks, panel peramalan interaktif, alur terpandu (\textit{guided flow}), serta manajemen riwayat percakapan.
\end{enumerate}

\section{Saran}

Berdasarkan hasil evaluasi dan keterbatasan yang diidentifikasi selama pelaksanaan tugas akhir, berikut adalah saran untuk pengembangan sistem di masa mendatang.

\begin{enumerate}
  \item \textbf{Integrasi Model Bahasa Besar (\textit{Large Language Model}/LLM)}: Pendekatan klasifikasi maksud berbasis pencocokan pola kata kunci memiliki keterbatasan dalam menangani variasi bahasa alami yang kompleks. Integrasi dengan LLM seperti GPT atau Gemini dapat meningkatkan kemampuan pemahaman bahasa alami secara signifikan, khususnya untuk menangani kueri ambigu, parafrase, dan pertanyaan di luar cakupan aturan yang telah didefinisikan.

  \item \textbf{Implementasi Model \textit{Temporal Fusion Transformer} (TFT)}: Model TFT yang telah dirancang pada tahap perancangan belum diimplementasikan secara penuh. Implementasi TFT menggunakan pustaka \texttt{pytorch-forecasting} diharapkan dapat meningkatkan akurasi peramalan melalui mekanisme perhatian multi-kepala (\textit{multi-head attention}) dan kemampuan menangkap pola temporal jangka panjang.

  \item \textbf{Integrasi dengan Sistem ERP/CRM Perusahaan}: Sistem yang dikembangkan saat ini merupakan purwarupa mandiri (\textit{standalone}). Integrasi dengan sistem ERP (\textit{Enterprise Resource Planning}) atau CRM (\textit{Customer Relationship Management}) yang ada di perusahaan dapat memperluas cakupan data dan fungsionalitas, serta mendukung alur kerja bisnis yang lebih komprehensif.

  \item \textbf{Peningkatan Keamanan ke Standar \textit{Enterprise}}: Implementasi keamanan saat ini berada pada tingkat dasar (autentikasi JWT dan enkripsi \textit{bcrypt}). Peningkatan ke standar keamanan \textit{enterprise} seperti ISO 27001, implementasi \textit{OAuth 2.0}, \textit{audit logging}, dan enkripsi data \textit{at-rest} diperlukan untuk penerapan pada lingkungan produksi.

  \item \textbf{Perluasan Cakupan \textit{Intent} dan Domain Data}: Cakupan \textit{intent} saat ini dibatasi pada tujuh kategori utama. Perluasan cakupan \textit{intent} dapat dilakukan melalui analisis log percakapan pengguna secara berkala untuk mengidentifikasi pola pertanyaan baru yang belum tercakup. Selain itu, penambahan domain data baru di luar lima domain yang ada dapat meningkatkan nilai analitik sistem.

  \item \textbf{Kontainerisasi dan \textit{Deployment} Otomatis}: Implementasi kontainerisasi menggunakan Docker dan orkestrasi dengan Kubernetes atau Docker Compose dapat meningkatkan konsistensi lingkungan \textit{deployment}, skalabilitas, dan kemudahan pemeliharaan sistem pada lingkungan produksi.

  \item \textbf{Perluasan Validasi ke Kelompok Pengguna yang Lebih Luas}: Validasi sistem saat ini telah dilakukan oleh pengguna internal dari tim \textit{Data Management}. Ke depannya, validasi perlu diperluas kepada kelompok pengguna yang lebih beragam di luar tim \textit{Data Management}, mencakup pengguna dari fungsi bisnis lain seperti tim penjualan, tim perencanaan, dan manajemen operasional. Perluasan ini bertujuan untuk mengukur efektivitas sistem dalam mendukung pengambilan keputusan pada berbagai konteks bisnis, mengidentifikasi kebutuhan analitik baru yang belum tercakup, serta memperoleh umpan balik yang lebih representatif mengenai kemudahan penggunaan dan relevansi respons sistem.
\end{enumerate}

